\section{Inner Product Spaces}

	One of the structures that we can put on a vector space is an inner product.  

  \begin{definition}[Inner Product]
		An \textbf{inner product} on a vector space $V$ over field $\mathbb{F} = \mathbb{R}$ or $\mathbb{C}$ is a mapping 
    \begin{equation}
      \langle \cdot, \cdot \rangle: V \times V \longrightarrow \mathbb{F}
    \end{equation}
    satisfying three properties 
    \begin{enumerate}
			\item \textit{First Argument Linearity}.\footnote{When the field is $\mathbb{C}$, the inner product is \textbf{sesqui-linear}, that is, linear with respect to the first argument and \textbf{skew linear} with respect to the second. When $\mathbb{R}$, it is bilinear.} $\langle \lambda x + \mu y, z \rangle = \lambda \langle x, z \rangle + \mu \langle y, z \rangle$
			\item \textit{Conjugate symmetry}. $\langle x, y \rangle = \overline{\langle y, x \rangle}$
      \item \textit{Positive Definiteness}. $\langle x, x \rangle \geq 0$, with $\langle x, x \rangle = 0 \iff x = 0$
    \end{enumerate}
		A vector space $V$ with an inner product is called an \textbf{inner product space}. 
  \end{definition}

	An inner product allows us to define some notion of an angle between two vectors in $V$. The inner product of a vector space $V$ over $\mathbb{R}$ is an element of $V^* \otimes V^*$. This concept of the metric tensor occurs when studying Riemannian manifolds in general relativity. 

  \begin{theorem}[Cauchy-Schwarz Inequality]
		Let $V$ be a vector space over $\mathbb{C}$. For all $v, w \in V$,\footnote{This inequality is often written $| \langle v, w \rangle | \leq \|v\| \|w\|$, but we must first prove that the norm induced by an inner product is actually norm. }
    \begin{equation}
			| \langle v, w \rangle | \leq \sqrt{\langle v, v \rangle} \sqrt{\langle w, w \rangle}
    \end{equation}
  \end{theorem}
	\begin{proof}
		We use a variational method. First, if $\|w\| = 0$, then $w = 0$ so by first argument linearity, $\langle v, w \rangle = 0$ and the inequality trivially holds. Now for the case where $w \neq 0 $, note that by positive definiteness of the norm, for all $t \in \mathbb{C}$, we have 
		\begin{align}
			0 & \leq \|v - tw\|^2 \\ 
				& = \langle v - tw, v - tw \rangle \\ 
				& = \langle v, v \rangle - \langle v, tw \rangle - \langle tw, v \rangle + \langle tw, tw \rangle \\ 
				& = \|v\|^2 - \overline{t} \langle v, w \rangle - t \overline{\langle v, w \rangle} + t \overline{t} \|w\|^2 \\ 
				& = \|v\|^2 - 2 \mathrm{Re}\{\overline{t} \langle v, w \rangle\} + t \overline{t} \|w\|^2
		\end{align}
		Therefore, by setting 
		\begin{equation}
			t = \frac{\langle v, w \rangle}{\|w\|^2}, 
		\end{equation}
		we get 
		\begin{align}
			0 & \leq \|v\|^2 - 2 \frac{\mathrm{Re}\{\langle v, w \rangle \overline{\langle v, w \rangle}\}}{\|w\|^2} + \frac{\langle v, w \rangle \overline{\langle v, w \rangle}}{\|w\|^4} \|w\|^2 \\
				& = \|v\|^2 - 2 \frac{\mathrm{Re}\{|\langle v, w \rangle|^2\}}{\|w\|^2} + \frac{|\langle v, w \rangle|^2}{\|w\|^2} \\
				& = \|v\|^2 - 2 \frac{|\langle v, w \rangle|^2}{\|w\|^2} + \frac{|\langle v, w \rangle|^2}{\|w\|^2} \\
				& = \|v\|^2 - \frac{|\langle v, w \rangle|^2}{\|w\|^2}
		\end{align}
		Since $|\langle v, w \rangle|$ is a real number. So, rearranging gives $|\langle v, w \rangle|^2 \leq \|v\|^2 \|w\|^2$, and square rooting both sides gives the inequality. 
	\end{proof}
	\begin{proof}
		For our second proof, we can restrict ourselves to real vector spaces and apply the same method to get 
		we apply the steps in the previous proof to get  
		\begin{equation}
		  0 \leq \|v\|^2 - 2 t \langle v, w \rangle + t^2 \|w\|^2
		\end{equation} 
		This is a quadratic form with respect to $t$, and since it is always greater than $0$, it has no solutions. So the discriminant must be less than $0$. 
		\begin{equation}
			(2 \langle v, w \rangle)^2 - 4 \|w\|^2 \|v\|^2 < 0
		\end{equation}
		Rearranging gives $|\langle v, w \rangle|^2 \leq \|v\|^2 \|w\|^2$, and square rooting both sides gives the inequality. 
	\end{proof}

  \begin{example}[Dot Product]
    Given vectors $v, w \in \mathbb{R}^n$, 
    \begin{equation}
      v \cdot w = \begin{pmatrix} v_1 \\ v_2 \\ \vdots \\ v_n \end{pmatrix} \cdot \begin{pmatrix} w_1 \\ w_2 \\ \vdots \\ w_n \end{pmatrix} \coloneqq \sum_{i=1}^n v_i w_i
    \end{equation}
  \end{example}

  \begin{example}[Integral Product]
    We haven't established any notion of integration yet, but after learning analysis, we can say the following. Let $C^0[a, b]$ be the space of all continuous real-valued functions defined over the interval $[a,b] \subset \mathbb{R}$. Given $f, g \in C^0[a,b]$, 
    \begin{equation}
      \langle f, g \rangle \coloneqq \int_a^b f(x)g(x) d x
    \end{equation}
    is an inner product on $C^0[a, b]$. 
  \end{example}

\subsection{Induced Norms}

  \begin{lemma}[Inner Product Induces Norm]
    An inner product induces a norm in the following way
    \begin{equation}
      \|v\| \coloneqq \sqrt{\langle v, v \rangle} 
    \end{equation}
  \end{lemma}
	\begin{proof}
		We must prove that the following norm satisfies the three properties of a norm: 
		\begin{enumerate}
			\item \textit{Positive Definite}. For any vector $v \in V$, 
				\begin{equation}
					\|v\| = \sqrt{\langle v, v \rangle} = 0 \implies \langle v, v \rangle = 0 \implies v = 0
				\end{equation}
			\item \textit{Absolute Homogeneity}. For any vector $v \in V$ and $c \in \mathbb{F}$, 
				\begin{equation}
					\|c v\| = \sqrt{\langle c v, cv \rangle} = \sqrt{c^2 \langle v, v \rangle} = |c| \sqrt{\langle v, v \rangle} = |c| \|v\|
				\end{equation}
		  \item \textit{Triangle Inequality}. For all $v, w \in V$, we can use Cauchy-Schwartz to get
				\begin{align}
					\|v + w\|^2 & = \|v\|^2 + \langle v, w \rangle + \langle w, v \rangle + \|w\|^2 \\  
											& \leq \|v\|^2 + 2 \|v\|\|w\|  + \|w\|^2 \\  
											& = (\|v\| + \|w\|)^2 
				\end{align}
				Then square rooting both sides gives $\|v + w\| \leq \|v\| + \|w\|$. 
		\end{enumerate}
	\end{proof}

	Now combining the following theorems give the fact that if $\| \cdot \|$ is the norm induced by an inner product $\langle \cdot , \cdot \rangle$, then $|\langle v, w \rangle| \leq \|v\| \|w\|$. 

	\begin{theorem}[Parallelogram Identity]
		Let $V$ be a real vector space with norm $\|v\| \coloneqq \sqrt{\langle v, v \rangle}$. Then, the identity holds for all $u, v \in V$. 
		\begin{equation}
		  \|u + v\|^2 + \|u - v\|^2 = 2 \|u\|^2 + 2 \|v\|^2
		\end{equation}
	\end{theorem}
	\begin{proof}
		We simply expand 
		\begin{align}
			\|u + v\|^2 + \|u - v\|^2 & = \langle u, u \rangle + 2 \langle u, v \rangle + \langle v, v \rangle + \langle u, u \rangle - 2 \langle u, v \rangle + \langle v, v \rangle \\ 
																& = 2 \langle u, u \rangle + 2 \langle v, v \rangle \\ 
																& = 2 \|u\|^2 + 2 \|v\|^2
		\end{align}
	\end{proof}

	Similar to \hyperref[pst-def:metrizable_space]{metrization of topological spaces}, a natural question to ask is whether all norms are induced by some inner product. The answer is no. 
	
	\begin{example}[L1 Norm]
		The L1 norm on $\mathbb{R}^n$ is not induced by any inner product, since if it was, the parallelogram identity will hold for all $u, v \in \mathbb{R}^2$. But for $u = (1, 0), v = (0, 1)$, and $u + v = (1, 1), u - v = (1, -1)$, we have 
		\begin{equation}
		  \|u + v\|^2 + \|u - v\|^2 = 2^2 + 2^2 = 8 \neq 4 = 2 + 2 = 2\|u\|^2 + 2 \|v\|^2
		\end{equation}
	\end{example}

	\begin{theorem}[Variational Definition of Norm]
    \begin{equation}
      \|v\| = \max_{\|w\|=1} \langle v, w \rangle
    \end{equation}
  \end{theorem} 
	\begin{proof}
	  It suffices to prove the following inequalities. 
		\begin{enumerate}
			\item $\max_{\|w\| = 1} \langle v, w \rangle \leq \|v\|$. Let $v \in V$ and $w \in V$ be a unit vector. Then, 
				\begin{equation}
				  \langle v, w \rangle \leq \|v\| \|w\| = \|v\|
				\end{equation}
				Therefore, by taking the maximum of both sides we get the desired result. 
			\item $\max_{\|w\| = 1} \langle v, w \rangle \geq \|v\|$. Consider the case when $w = \frac{v}{\|v\|}$. Then, 
				\begin{equation}
					\langle v, w \rangle = \langle v, \frac{v}{\|v\|} \rangle = \frac{1}{\|v\|} \langle v, v \rangle = \frac{\|v\|^2}{\|v\|} = \|v\|
				\end{equation}
				Therefore, the maximum of $\langle v, w \rangle$ must be at least $\|v\|$. 
		\end{enumerate}
	\end{proof}

\subsection{Orthogonality}

  \begin{definition}[Orthogonal Vectors]
    Two vectors $v, w \in V$ of an inner product space are said to be \textbf{orthogonal} if 
    \begin{equation}
      \langle v, w \rangle = 0
    \end{equation}
  \end{definition}

  Note that the definition of orthogonality is dependent on the definition of the inner product. If the inner product is defined differently, then orthogonality will be defined differently. In the case when the inner product is defined to be the dot product, orthogonality is defined to be the "normal" perpendicularity between vectors. We can further define subspaces to be orthogonal. 

  \begin{definition}[Orthogonal Subspaces]
    Two subspaces $U, W$ of inner product space $V$ are said to be orthogonal to each other if 
    \begin{equation}
      \langle u, w \rangle = 0, \qquad  u \in U, w \in W
    \end{equation}
  \end{definition}

  \begin{definition}[Orthogonal Complement]
    Given a subspace $W$ of inner product space $V$, the \textbf{orthogonal complement} of $W$, denoted $W^\perp$, is defined
    \begin{equation}
      \{ v \in V \colon \langle v, w \rangle = 0 \quad  \forall w \in W\}
    \end{equation}
    which is the set of all vectors in $V$ orthogonal to every vector in $W$. 
  \end{definition}

	\begin{definition}[Orthogonal, Orthonormal Basis]
		A basis $\{v_i\}_i$ of an inner product space $V$ is said to be \textbf{orthogonal} if for every pair $i, j$ where $i \neq j$, $\langle v_i, v_j \rangle = 0$. It is said to be \textbf{orthonormal} if 
		\begin{equation}
			\langle v_i, v_j \rangle = \begin{cases} 
				1 & \text{ if } i = j \\
 				0 & \text{ if } i \neq j 
			\end{cases}\footnote{This is often called the \textit{Kronecker delta function}.}
		\end{equation}
	\end{definition}

\subsection{Projections}

  \begin{definition}[Projection]
    A \textbf{projection mapping} is a linear mapping $P$ where 
    \begin{equation}
      P = P^2
    \end{equation}
  \end{definition}

  The concept of orthogonality also allows us to define orthogonal projections onto a vector or subspace. 

	\begin{lemma}[Orthogonal Decomposition Theorem]
		Given subspace $W$ of inner product space $V$, 
	  \begin{equation}
			 W \oplus W^\perp = V  
	  \end{equation}
	\end{lemma}
	\begin{proof}
		Take any vector $v \in V$, and take
		\begin{equation}
			w = \argmin_{u \in W} \|u - v\|^2
		\end{equation}
		which is well defined since subspace $W$ is \hyperref[pst-def:closed]{closed}, $u \to \|u - v\|^2$ is \hyperref[pst-def:continuous_func]{continuous}, and so by the \hyperref[pst-thm:evt]{Extreme Value Theorem}, the minimum is attained. Now set 
		\begin{equation}
			w^\perp = v - w. 
		\end{equation}
		We wish to show that $\langle w^\perp, u \rangle = 0$ for all $u \in W$. For sake of contradiction, assume that this isn't true, and we want to arrive at a result that contradicts the fact that $w$ is really the minimizer. So, there exists some $u_0 \in W$ s.t. $\langle w^\perp, u_0 \rangle = c \neq 0$. WLOG assume $\|u_0\| = 1$. Choose a new vector 
		\begin{equation}
			u_1 = w + \overline{c} u_0 \in W
		\end{equation}
		where $\overline{c}$ is the complex conjugate of $c$. Note that $u_1$---as a linear combination of vectors in $W$---is also in $W$. Now let's take a look at the distance between $u_1$ and $v$. 
		\begin{align}
			\|v - u_1\|^2 & = \|v - w - \overline{c} u_0\|^2 \\ 
										& = \|w^\perp - \overline{c} u_0\|^2 \\ 
										& = \langle w^\perp - \overline{c} u_0, w^\perp - \overline{c} u_0 \rangle \\ 
										& = \langle w^\perp, w^\perp \rangle - \overline{c} \langle w^\perp, u_0 \rangle - \overline{c} \langle u_0, w^\perp \rangle + c \overline{c} \langle u_0, u_0 \rangle \\ 
										& = \|w^\perp\|^2 - \overline{c} c - c \overline{c} + c \overline{c} \\ 
										& = \|w^\perp\|^2 - |c|^2 < \|w^\perp\|^2
		\end{align}
		This contradicts the fact that $w$ minimizes $u \mapsto \|u - v\|^2$, and we are done. Now to prove uniqueness, assume that there is another decomposition $v = w^\prime + w^{\prime\perp}$. Then, 
		\begin{equation}
			w + w^\perp = w^\prime + w^{\prime \perp} \implies w - w^\prime = w^{\prime\perp} - w^\prime 
		\end{equation}
		But the LHS is a vector in $W$ and the RHS is a vector in $W^\perp$. The only such vector in both subspaces is $0$. 
	\end{proof}

	With this decomposition theorem, the orthogonal projection operator is well-defined. 

  \begin{definition}[Orthogonal Projection]
    Let $v \in V$ and let $W$ be a subspace of $V$. The \textbf{orthogonal projection} of $v$ onto $W$ is then defined as 
    \begin{equation}
      \text{proj}_W (v) = w
    \end{equation}
		where $v = w + w^\perp$, where $w \in W, w^\perp \in W^\perp$, is the orthogonal decomposition of $v$. 
  \end{definition}

	\begin{lemma}[Orthogonal Projections are Linear]
    Orthogonal projections are linear transformations. 
	\end{lemma}
	\begin{proof}
		Let $u, v \in V$, $c \in \mathbb{F}$, and $W$ be a subspace of $V$. Then, the orthogonal decompositions are $u = u_W + u_W^\perp$, $v = v_W + v_W^\perp$. Therefore, the following are valid orthogonal decompositions of $u + v$ and $cu$: 
		\begin{align}
			u + v & = (u_W + v_W) + (u_W^\perp + v_W^\perp) \\
			c u & = cu_W + cu_W^\perp 
		\end{align}
		Since orthogonal decompositions are unique, this defines the linear map. 
	\end{proof}

	\begin{lemma}[Orthogonal Projection Formula onto 1-Dimension in $\mathbb{R}^n$]
		Let $W = \mathrm{span}\{w\}$ be a $1$-dimensional subspace of $\mathbb{R}^n$. Then, the orthogonal projection of $v$ onto $W$ is 
		\begin{equation}
			\mathrm{proj}_W (v) = \frac{\langle v, w \rangle}{\|w\|^2} w
		\end{equation}
	\end{lemma}
	\begin{proof}
		By definition of the orthogonal projection, we are looking for the minimizer 
		\begin{equation}
			\argmin_{u \in W} \|v - u\|^2 
		\end{equation}
		Since $u \in W$, we can simply rewrite this problem as 
		\begin{equation}
			\argmin_{t \in \mathbb{R}} \|v - tw\|^2 
		\end{equation}
		By expanding, we need to minimize the quadratic form 
		\begin{equation}
		  \langle v, v \rangle - 2 t \langle v, w \rangle + t^2 \langle w, w \rangle
		\end{equation} 
		with respect to $t$. This is a parabola that clearly opens up, and so the minimizer is 
		\begin{equation}
			t = - \frac{-2 \langle v, w \rangle}{2 \langle w, w \rangle} = \frac{\langle v, w \rangle}{\|w\|^2}
		\end{equation}
	\end{proof}

	\begin{theorem}[Orthogonal Projection Formula in $\mathbb{R}^n$]
		Given $k$-dimensional subspace $W$ of $\mathbb{R}^n$ with an \textit{orthogonal} basis $\{w_1, w_2, \ldots, w_k\}$, the orthogonal projection of $v \in \mathbb{R}^n$ to $W$ can be computed with the formula 
    \begin{equation}
			\text{proj}_W (v) = \sum_{i=1}^k \text{proj}_{w_i} (v) = \sum_{i=1}^k \frac{\langle v, w_i \rangle}{\|w_i\|^2} w_i 
    \end{equation}
  \end{theorem}
	\begin{proof}
		Let's decompose $v$ into the basis of $w_i$'s, plus some orthogonal term in $W^\perp$.\footnote{I first tried to prove it by letting $u = \proj_W (v) \implies v = w + w^\perp$, and $u_i = \proj_{w_i} (u) \implies v = u_i + u_i^\perp$. Then I tried setting $w + w^\perp = u_i + u_i^\perp$ for all $i$, but this requires me to prove some properties of the projection with respect to the vector being projected on. So I'm not allowed to use linearity here and was stuck. }
		\begin{equation}
		  v = a_1 w_1 + \ldots + a_k w_k + w^\perp, \qquad w^\perp = W^\perp
		\end{equation} 
		By linearity of projection, 
		\begin{align}
			\proj_W (v) & = \bigg( \sum_{i=1}^k a_i \proj_W (w_i) \bigg) + \proj_W (w^\perp) \\ 
									& = \bigg( \sum_{i=1}^k a_i w_i \bigg) + 0 \\ 
									& = \sum_{i=1}^k a_i \proj_{w_i} (v)
		\end{align}
		Note that in the last line, we need the orthogonality of the $w_i$'s. 
	\end{proof}

  Given that we have an orthonormal basis $\{r_i\}_{i=1}^k$ of subspace $Y$ in $\mathbb{R}^n$, we can more simply define 
  \begin{equation}
    \text{proj}_Y (x) = \sum_{i=1}^k \langle x, r_i \rangle r_i
  \end{equation}

  \begin{theorem}[Gram-Schmidt]
    Every basis in a finite-dimensional vector space can be transformed into an orthonormal basis. 
  \end{theorem}
  \begin{proof}
		We start off with any basis, not necessarily orthonormal, of $X$, denoted $\{x_1, x_2, \cdots, x_n\}$. We first assign 
    \begin{equation}
      x_1 = l_1
    \end{equation}
    Then we take $x_2$ and find the orthogonal component (with respect to $l_1$) with the equation
    \begin{equation}
      l_2 = x_2 - \text{proj}_{l_1} (x_2)
    \end{equation}
    This creates an orthogonal basis for span$\{x_1, x_2\}$. Then we take $x_3$ and find the orthogonal component (with respect to span$\{l_1, l_2\}$. 
    \begin{equation}
      l_3 = x_3 - \text{proj}_{l_1} (x_3) - \text{proj}_{l_2} (x_3)
    \end{equation}
    This creates an orthogonal basis for span$\{x_1, x_2, x_3\}$. We repeat this process until we complete the basis of $X$, using the general equation
    \begin{equation}
      l_k = x_k - \sum_{i=1}^{k-1} \text{proj}_{l_k} (x_k) = x_k - \sum_{i=1}^{k-1} \frac{x_k \cdot l_k}{||l_k||^2} l_k, \; k = 1, 2, \cdots, n
    \end{equation}
    Finally, we take these orthogonal vectors and normalize them to magnitude 1. Note that this algorithm does not produce a unique orthonormal basis. Rather, it is highly un-unique. 
  \end{proof}

\subsection{Isometries}

  \begin{definition}[Isometry]
    An \textbf{isometry} $M$ of metric space $(X, d)$ is a mapping that preserves all distances. That is, for all $x, y \in X$, 
    \begin{equation}
      d(x, y) = d(M x, M y) 
    \end{equation}
    The set of all isometries, denoted Isom$(X)$, is a group that is generated by all translations, rotations, and reflections. 
  \end{definition}

  \begin{example}
    A rotation around any axis or a flip across any hyperplane are all isometries. 
  \end{example}

  Since linear maps always preserve the origin, we will focus on origin-preserving isometries, which is a subgroup called the orthogonal group.

  \begin{definition}[Orthogonal Group]
		The \textbf{orthogonal group}, denoted $O(n)$, is the group of all linear isometries of $\mathbb{R}^n$. 
  \end{definition}

  \begin{definition}[Special Orthogonal Group]
		The \textbf{special orthogonal group}, denoted $SO(n)$, is the group of all isometries of $\mathbb{R}^n$ that preserve the handedness. 
  \end{definition}

	$SO(n)$ is a subgroup of $O(n)$. We extend this concept to complex Euclidean spaces. 

  \begin{definition}[Unitary Group]
		The \textbf{unitary group}, denoted $U(n)$ is the group of all isometrices of $\mathbb{C}^n$. 
  \end{definition}

  \begin{example}
    $U(1)$ is the set of complex numbers with norm $1$. 
  \end{example}

  \begin{definition}[Special Unitary Group]
		The \textbf{special unitary group}, denoted $SU(n)$ is the group of all isometrices of $\mathbb{C}^n$ that preserve the handedness. 
  \end{definition}

  The groups mentioned in this section are examples of \textit{Lie Groups}. Lie groups in general will not be defined in here, since they require knowledge of smooth manifolds and differential geometry. In order to analyze these abstract groups, we use the exponential map $e \in$ End$($ Mat$(n, \mathbb{F})$ to reduce these Lie groups to Lie algebras.

\subsection{Adjoint Operators}

	Let $A: U \to V$ be a linear mapping between inner product spaces, with the inner product in $U$ and $V$ denoted $\langle \cdot,\cdot \rangle_U$ and $\langle \cdot,\cdot \rangle_V$, respectively. First, fix any $v \in V$ and define the linear function $l \in U^\ast$
	\begin{equation}
		l(\cdot) \coloneqq \langle A(\cdot), v \rangle_V
	\end{equation}
	Since $U$ is naturally isomorphic to $U^\ast$, for this specific $l$, we can find a $u \in U$ such that
	\begin{equation}
		l(\cdot) \coloneqq \langle A(\cdot), v \rangle_V = \langle \cdot, u \rangle_U
	\end{equation} 
	Such a map from $v \mapsto u$ is linear (should verify this), and we call this map the \textit{adjoint operator} $A^\dagger: V \to U$. Therefore, we can write the equation above to 
	\begin{equation}
	  \langle A(\cdot), v \rangle_V = \langle \cdot, A^\dagger (u) \rangle
	\end{equation}

  \begin{definition}[Adjoint Operator]
		Let $A: U \to V$ be a linear mapping between inner product spaces. Then, the \textbf{adjoint operator} of $A$ is any operator $A^\dagger: V \to U$ such that
		\begin{equation}
		  \langle u, A^\dagger v \rangle_U = \langle A u, v \rangle_V
		\end{equation}
		for all $u, v \in U$. 

		\begin{figure}[H]
			\centering 
			\begin{tikzcd}
				U \arrow{r}{A} \arrow{d}{i} & V \arrow{d}{j}\\
				U^\ast & V^\ast \arrow{l}{A^T}
			\end{tikzcd}
			\caption{The adjoint operator is the map $A^\dagger$ such that the diagram commutes.} 
			\label{fig:adjoint}
		\end{figure}
  \end{definition}

  It is important to note that the adjoint is not the same as the transpose since the transpose is a mapping between the dual spaces. Furthermore, the transpose is canonically defined upon defining the linear transformation $A: U \to V$, while defining the adjoint requires the additional structure of an isomorphism from $U$ to $U^\ast$ and from $V$ to $V^\ast$. There are two ways to define these isomorphisms. 

  First, we can define dot products on both $U$ and $V$ and define the natural isomorphism 
  \begin{align*}
    &i: U \longrightarrow U^\ast, \; i(u) \equiv (u, \cdot) \in U^\ast\\
    &j: V \longrightarrow V^\ast, \; j(v) \equiv (v, \cdot) \in V^\ast
  \end{align*}
  This canonically creates the mapping 
  \begin{equation}
    i^{-1} A^T j: V \longrightarrow U
  \end{equation}
  which we define as the adjoint $A^\dagger$. This method using natural isomorphisms is precisely how we have defined the adjoint above. There is a second way, however. We can fix \textbf{orthonomal} bases on $U$ and $V$ and then assign them their respective dual spaces (satisfying the Kronecker delta function). Let the basis of $U$ be $\{u_1, \cdots, u_n\}$, $U^\ast$ be $\{u_1^\prime, \cdots, u_n^\prime\}$, $V$ be $\{v_1, \cdots, v_m\}$, and $V^\ast$ be $\{v_1^\prime, \cdots, v_m^\prime\}$. Now we can define the isomorphisms 
  \begin{align*}
    & i^\prime: U \longrightarrow U^\ast, \; i^\prime (u) \equiv c_1 u_1^\prime + \cdots + c_n u_n^\prime \\
    & j^\prime: V \longrightarrow V^\ast, \; j^\prime (v) \equiv k_1 v_1^\prime + \cdots + k_m v_m^\prime
  \end{align*}
  and then define the adjoint as 
  \begin{equation}
    A^\dagger \equiv i^{\prime -1} A^T j^\prime
  \end{equation}
  Let us compare these two definitions. Given a vector $u = a_1 u_1 + \cdots + a_n u_n, \Tilde{u} = b_1 u_1 + \cdots b_n u_n \in U$, 
  \begin{align*}
    &i(u)(\Tilde{u}) \equiv (u, \Tilde{u}) = \Big(\sum_{\alpha = 1}^n a_\alpha u_\alpha, \sum_{\beta = 1}^n b_\beta u_\beta \Big) = \sum_{\alpha, \beta} a_\alpha b_\beta \delta^\alpha_\beta = \sum_{\gamma=1}^n a_\gamma b_\gamma \\
    &i^\prime(u)(\Tilde{u}) \equiv \Big(\sum_{i=1}^n a_i u_i^\prime \Big) \Big( \sum_{j=1}^n b_j u_j \Big) = \sum_{i, j} a_i b_j u_i^\prime (u_j) = \sum_{i, j} a_i b_j \delta^i_j = \sum_{k=1}^n a_k b_k
  \end{align*}
  Similarly for vector $v = g_1 v_1 + \cdots g_n v_n, \Tilde{v} = h_1 v_1 + \cdots + h_n v_n \in V$, 
  \begin{align*}
    &i(v)(\Tilde{v}) \equiv (v, \Tilde{v}) = \Big(\sum_{\alpha = 1}^n g_\alpha v_\alpha, \sum_{\beta = 1}^n h_\beta v_\beta \Big) = \sum_{\alpha, \beta} g_\alpha h_\beta \delta^\alpha_\beta = \sum_{\gamma=1}^n g_\gamma h_\gamma \\
    &i^\prime(v)(\Tilde{v}) \equiv \Big(\sum_{i=1}^n g_i u_i^\prime \Big) \Big( \sum_{j=1}^n h_j v_j \Big) = \sum_{i, j} g_i h_j v_i^\prime (v_j) = \sum_{i, j} g_i h_j \delta^i_j = \sum_{k=1}^n g_k h_k
  \end{align*}
  Therefore, $i = i^\prime$ and $j = j^\prime$, meaning that the two derivations of the adjoint $A = i^{-1} A^T j = i^{-1 \prime} A^T j^\prime$ are exactly the same! We must note that the basis endowed on both $U$ and $V$ must be orthonormal for it to "mimic" the inner product. The derivation of the adjoint in these two equivalent methods may help the reader further understand that the adjoint $A^\dagger$ is really just a composition of fundamental linear functions $j: V \longrightarrow V^\ast$, $A^T: V^\ast \longrightarrow U^\ast$, and $i^{-1}: U^\ast \longrightarrow U$ that are all canonically created as soon as $A: U \longrightarrow V$ is created, along with the inner product spaces $U$ and $V$. 

  However, it is hard to grasp a visual intuition of adjoint operators in general. Note that the properties of the transpose indicate that given $A: \mathbb{R}^n \longrightarrow \mathbb{R}^m$ with the standard orthonormal basis and dot product, the matrix representation of $A^\dagger$ is just $A^T$. If $A$ is a matrix over $\mathbb{C}$, then $A^\dagger$ is $A^H \equiv \bar{A}^T$, the \textbf{Hermitian transpose}, or \textbf{conjugate transpose}, of $A$. 

  Note that this definition of the adjoint of linear operators is completely unrelated to the definition of an adjoint of a matrix! 

  We now describe one common application of adjoints. 
  \begin{theorem}
    Let $A \in$ Mat$(m \times n, \mathbb{R})$ with $m > n$. This means that the system of equations $A x = p$ is an overdetermined system and will have no solutions with probability 1. However, we can find the \textbf{best-fit solution} of the system. That is, the vector $x$ that minimizes $\|A x -p\|^2$ is the solution $z$ of 
    \begin{equation}
      A^\dagger A z = A^\dagger p
    \end{equation}
    $z$ is therefore, the "closest approximation" of the solution of $A x = p$ that lives in $\mathbb{R}^n$. 
  \end{theorem}

  \begin{theorem}[Properties of the Adjoint] 
		Let $A, B: X \longrightarrow U, \; C: U \longrightarrow V$ be linear mappings. Then, 
    \begin{enumerate}
      \item $(A + B)^\dagger = A^\dagger + B^\dagger$
      \item $(C A)^\dagger = A^\dagger C^\dagger$
      \item $(A^{-1})^\dagger = (A^\dagger)^{-1}$ if $A$ is bijective
      \item $(A^\dagger)^\dagger = A$
    \end{enumerate}
  \end{theorem}

	\begin{definition}[Self-Adjoint Map]
    Linear mapping $A$ is \textbf{self-adjoint} if and only if $A = A^\dagger$. 
  \end{definition}

	If $M$ is any linear mapping, then its self-adjoint part is 
	\begin{equation}
		M_\delta = \frac{M + M^\dagger}{2}
	\end{equation}

	\begin{theorem}[Orthogonal Projections are Self-Adjoint]
		Let $P$ be an orthogonal projection. Then, $P = P^\dagger$. 
  \end{theorem}


  \begin{definition}[Anti-Self-Adjoint]
    Map $A$ is \textbf{anti-self adjoint} if $A^\dagger = - A$. Conjugate symmetry implies that
    \begin{equation}
      A^\dagger = A \iff (i A)^\dagger = - (i A)
    \end{equation}
  \end{definition}

\subsection{Dual Inner Product Spaces}

	\begin{theorem}[Basis-Free Isomorphism Induced by Inner Product]
    The inner product endowed on $V$ induces a natural isomorphism between $V$ and $V^\ast$. 
  \end{theorem}
  \begin{proof}
    We fix $y \in V$ and simply define the isomorphism to be. 
    \begin{equation}
      l(y) \equiv (x, y)
    \end{equation}
    which defines a bijection between $x \in V$ and $l \in V^*$. 
  \end{proof}

  Note that given vector spaces $U, V$, the set of all linear mappings $A: U \longrightarrow V$ also forms a vector space. More specifically, it is a rank (1,1) tensor product space. This means that we can define similar Euclidean structures on them. The norm of a matrix is worth mentioning. Note that the structures and concepts of metrics, norms, inner products, distances, magnitudes, orthogonality, and basis are not intrinsic properties of the vector space. So, we will not assume the existence of these structures unless otherwise stated or explicitly implied. 

